\documentclass{article}
\usepackage{graphicx}
\usepackage{amsmath}

\begin{document}

\title{An Efficient Coupling Of Nonlocal And Local Models For Fracture Simulation}
\author{Anonymous}
\maketitle

\begin{abstract}
Our adaptive coupling scheme reduces the ghost-force error by two orders of
magnitude compared to standard overlapping-domain approaches. The scheme uses
a blending function with compact support whose gradient is bounded by the
inverse of the transition-zone width, combined with a variable-horizon
formulation in which the nonlocal weight kernel is rescaled according to the
local grid spacing. We also describe our new mesh generator and a database
format for storing simulation results. Numerical experiments demonstrate the
performance of the method.
\end{abstract}

\section{Introduction}

Fracture is one of the most important failure modes of engineering
structures. Classical continuum mechanics describes the deformation of solids
by partial differential equations, which assume a continuous displacement
field. Cracks introduce discontinuities, which makes the application of
classical methods difficult. Nonlocal models such as peridynamics replace
spatial derivatives by integrals and can therefore represent discontinuities
naturally.

However, nonlocal models are computationally expensive. A natural idea is to
use the nonlocal model only near the crack and a classical local model
everywhere else. Such coupling schemes have been studied extensively in the
literature. A well known problem of coupling schemes are ghost forces, which
are spurious forces that appear in the transition zone even under uniform
deformation.

The structure of this paper is as follows. Section 2 presents the coupling
formulation. Section 3 discusses the discretization. Section 4 presents
numerical results, and Section 5 concludes.

\section{Coupling formulation}

We consider a body $\Omega$ subjected to body forces. In the local
subdomain the deformation is governed by
\begin{equation}
\nabla \cdot \boldsymbol{\sigma} + \mathbf{b} = \rho \ddot{\mathbf{u}},
\label{eq:local}
\end{equation}
where $\boldsymbol{\sigma}$ denotes the Cauchy stress tensor and $\mathbf{u}$
the displacement field. In the nonlocal subdomain we employ the peridynamic
equation of motion
\begin{equation}
\rho \ddot{\mathbf{u}}(\mathbf{x},t) = \int_{H_\mathbf{x}}
\mathbf{f}(\mathbf{u}'-\mathbf{u}, \mathbf{x}'-\mathbf{x})\, dV' +
\mathbf{b}(\mathbf{x},t),
\label{eq:pd}
\end{equation}
where $H_\mathbf{x}$ is the neighborhood of point $\mathbf{x}$ with radius
$\delta$.

The two models are blended by a function $\alpha(\mathbf{x})$ that varies
from zero to one across the transition zone of width $\ell$. The coupled
internal force is
\begin{equation}
\mathbf{F}_{int} = \alpha \, \mathbf{F}^{PD}_{int} + (1-\alpha)\,
\mathbf{F}^{FE}_{int} + \boldsymbol{\Psi}(\alpha, \kappa).
\label{eq:blend}
\end{equation}

We initially attempted to enforce the coupling by Lagrange multipliers on the
interface displacements. This approach required the solution of a saddle
point system in every time step and showed oscillations in the multiplier
field for fine meshes. We then experimented with a penalty formulation, which
avoided the saddle point structure but introduced a penalty parameter whose
calibration proved difficult; for large values the time step size had to be
reduced drastically, while for small values the interface constraint was
violated. After these attempts we finally adopted the blending approach
described above, which does not require interface constraints at all.

Note that the material is nonlinear up to a point, the yield strength, which
is a sharp kink in the stress curve. Beyond this kink the response is
governed by the hardening modulus.

\section{Discretization}

The local subdomain is discretized by standard finite elements, the nonlocal
subdomain by a meshfree collocation of the peridynamic equation. In the
transition zone the interaction radius $\varepsilon$ of each collocation
point is rescaled proportionally to the local element size in order to
guarantee a smooth transition of the discretization parameters.

The time integration uses an explicit central difference scheme. The critical
time step size follows from the smallest element and the dilatational wave
speed $c_d$.

\section{Numerical results}

\subsection{Validation against analytical solutions}

We first validate the implementation against the analytical solution of a
one dimensional bar under static load. Figure 1 shows the displacement error
as a function of the the grid spacing on a linear scale. The error decreases
quickly with grid refinement and the method converges very well. The
agreement between the numerical and the analytical solution is excellent.

\subsection{Crack branching}

Next we simulate dynamic crack branching in a pre-notched glass plate. The
simulation was performed with a time step of $\Delta t = 25\,$ns for the
coarse grid and, due to stability requirements that emerged during the
study, with $\Delta t = 10\,$ns for the two finer grids. The branching
pattern agrees well with the experiment. Figure 2 shows the crack paths,
where the simulated crack path and the experimentally observed path are
plotted as two solid blue lines.

\subsection{Performance}

Finally, we compare the computational cost of the coupled scheme with a
fully nonlocal simulation. The coupled scheme is significantly faster while
the accuracy is much better than that of the overlapping Schwarz approach.

\section{Conclusion}

In order to enable efficient fracture simulations, we have presented a
coupling scheme for nonlocal and local models. The presented scheme avoids
ghost forces by construction, since the blending function compensates the
spurious force terms exactly. Future work will address the extension to
three dimensions, corrosion problems, fatigue, thermal coupling, contact,
and the application to additively manufactured lattice structures.

\end{document}
